\texteparacol{\rubrique{\normaltext{Dóminus vobiscum,} etc. dicitur in die expositionis tantum.}

\vv Dóminus vo\textbf{bís}cum.

\rr Et cum spíritu \textbf{tu}o.

Orémus.

\lettrine[ante=1]{D}{e}us, qui nobis sub Sacraménto mirábili passiónis tuæ memóriam reliquísti:~† tríbue, quǽsumus, ita nos córporis, et sánguinis tui sacra mystéria venerári;~* ut redem\-ptiónis tuæ fru\-ctum in nobis júgiter sentiámus.

\lettrine[ante=2]{C}{o}ncéde nos famúlos tuos, quǽsumus Domíne Deus, perpétua mentis et corpóris sanitáte gaudére:~† et gloriósa beátæ Mariæ semper Virgínis intercessióne,~* a præsénti liberári tristítia, et æterna perfrúi lætítia.

\rubrique{Ab Adventu usque ad Nativitaem Domini:}

\lettrine{D}{e}us, qui de beátæ Maríæ Vírginis útero Verbum tuum, Angelo nuntiánte, carnem suscípere voluísti:~† præsta supplícibus tuis; ut, qui vere eam Genetrícem Dei credimus,~* ejus apud te intercessiónibus adjuvémur.

\rubrique{A Nativitate Domini usque ad Purificationem B.M.V.:} 

\lettrine{D}{e}us, qui salútis ætérnæ, beátæ Maríæ virginitáte fœcúnda, humáno géneri prǽmia præstitísti:~† tríbue, quǽsumus; ut ipsam pro nobis intercédere sentiámus,~* per quam merúimus auctórem vitæ suscípere, Dóminum nóstrum Jesum Christum Fílium tuum.

\rubrique{Tempore Paschali.}

\lettrine{D}{e}us, qui per resurrectiónem Fílii tui Dómini nostri Jesu Christi mundum lætificáre dignátus es:~† præsta, quǽsumus; ut per ejus Genetrícem Vírginem Maríam~* perpétuæ capiámus gáudia vitæ.

\lettrine[ante=3]{O}{m}nípotens sempitérne Deus, miserere fámulo tuo Pontífici
nostro \nomen:~† et dírige eum secúndum tuarn cleméntiam in
viam salútis ætérnae;~* ut te donánte, tibi plácita cúpiat, et tota
virtúte perfíciat.

\lettrine[ante=4]{D}{e}us, refúgium nostrum et virtus:~† adésto piis Ecclésiæ tuæ précibus, au\-ctor ipse pietátis, et præsta;~* ut quod fidéliter pétimus, efficáciter consequámur.

\lettrine[ante=5]{O}{m}nípotens sempitérne Deus, qui vivórum domináris simul et mortuórum, omniúmque miseréris, quos tuos fide et ópere futúros esse prænóscis:~† te súpplices exorámus; ut pro quibus effúndere preces decrévimus, quosque vel præsens sǽculum adhuc in carne rétinet vel futúrum jam exútos córpore suscépit,~* intercedéntibus ómnibus San\-ctis tuis, pietátis tuæ cleméntia, ómnium deli\-ctórum suórum véniam consequántur. Per Dóminum nostrum Jesum Christum.

\rr Amen.

\vv Dóminus vo\textbf{bís}cum.

\rr Et cum spíritu \textbf{tu}o.

\rubrique{Cantores.}

\vv Exáudiat nos omnípotens et miséricors Dóminus.

\rr Amen.

\rubrique{Celebrans, voce recta et depressa.}

\vv Et fidélium ánimæ per misericórdiam Dei requiéscant in pace.

\rr Amen.

\rubrique{In die expositionis, post brevem orationem, discedunt celebrans et sacri ministri.}

\rubrique{In die repositionis, sequitur Benedictio.}
}
{english}
{\rubrique{\normaltext{The Lord be with you,} etc. is said on the day of exposition only.}

\noindent	\vv The Lord be with you.
	
\noindent \rr And with thy spirit.

Let us pray.

\noindent 1.\@ O God, who in this wonderful Sacrament hast left us a memorial of thy Passion: grant, we implore thee, that we may so venerate the sacred mysteries of thy Body and Blood, as always to be conscious of the fruit of thy Redemption.

\noindent 2.\@ Grant, we beseech thee, o Lord God, unto us thy servants, that we may rejoice in continual health of mind and body; and, by the glorious intercession of Blessed Mary, ever Virgin, may we be delivered from present sadness and enter into the joy of thine eternal gladness.

\rubrique{From the first Sunday of Advent until Christmas Eve.}

\noindent O God, who, by the message of an Angel, didst will thy Word to take flesh in the womb of the Blessed Virgin Mary, grant that we thy suppliants, who believe her to be truly the Mother of God, may be helped by her intercession with thee.

\rubrique{From First Vespers of Christmas until the Purification.}

\noindent O God, who by the fruitful virginity of blessed Mary, hast given to the human race the rewards of eternal salvation: grant, we beseech thee, that we may experience her intercession for us, through whom we deserved to receive the Author of life, our Lord Jesus Christ, thy Son.

\rubrique{In Paschal Time.}

\noindent O God, who gave joy to the world through the resurrection of thy Son, our Lord Jesus Christ; grant, we beseech thee, that through his mother, the Virgin Mary, we may obtain the joys of everlasting life.

\noindent 3.\@ Almighty, everlasting God, have mercy upon thy servant, \nomen{} our Sovereign Pontiff, and direct him, according to thy clemency, into the way of everlasting salvation; that by thy grace he may desire those things that are pleasing to thee, and perform them with all his strength.

\noindent 4.\@ O God, our refuge and strength, who art the author of all godliness, be ready, we beseech thee, to hear the devout prayers of thy Church, and grant that those things which we ask faithfully, we may obtain effectually

\noindent 5. Almighty and everlasting God, who hast dominion over the living and the dead, and art merciful to all whom thou foreknowest shall be thine by faith and good works: we humbly beseech thee; that they for whom we intend to pour forth our prayers, whether this present world still detains them in the flesh, or the world to come has already received them out of their bodies, may, through the intercession of all thy saints, and in thy compassionate goodness, obtain the pardon of all their sins. Through Christ our Lord. Amen.

\noindent	\vv The Lord be with you.
	
\noindent \rr And with thy spirit.

\rubrique{Cantors.}

\noindent	\vv May the almighty and most merciful Lord graciously hear us.

\noindent \rr  Amen.

\rubrique{The celebrant, in a monotone, on a lower note.}

\noindent \vv May the souls of the faithful departed through the mercy of God rest in peace.

\noindent \rr Amen.

\rubrique{On the day of exposition, the celebrant and sacred ministers depart after a short prayer.}

\rubrique{On the day of reposition, Benediction follows.}}